% !TEX program = xelatex
% SG-Lean 思路文档（第二版）—— 2026-09-27 外部评审后重定稿
\documentclass[UTF8,a4paper,11pt,fontset=windows]{ctexart}

\usepackage{geometry}
\geometry{left=2.2cm,right=2.2cm,top=2.5cm,bottom=2.3cm}
\usepackage{amsmath,amssymb}
\usepackage{booktabs}
\usepackage{tabularx}
\usepackage{array}
\usepackage{enumitem}
\usepackage{xcolor}
\usepackage{url}
\usepackage[hidelinks]{hyperref}
\usepackage{fancyhdr}
\usepackage{tcolorbox}
\tcbuselibrary{skins,breakable}
\usepackage{tikz}
\usetikzlibrary{positioning,arrows.meta,fit}

\definecolor{sgblue}{RGB}{31,86,153}
\definecolor{sggreen}{RGB}{38,120,86}
\definecolor{sggray}{RGB}{90,96,106}
\definecolor{sgred}{RGB}{150,50,50}

\ctexset{
  section/format = \Large\bfseries\color{sgblue},
  subsection/format = \large\bfseries\color{sgblue!85!black},
  subsubsection/format = \normalsize\bfseries\color{sggray},
}

\pagestyle{fancy}
\fancyhf{}
\fancyhead[L]{\small\color{sggray}SG-Lean 实现规格}
\fancyhead[R]{\small\color{sggray}\leftmark}
\fancyfoot[C]{\small\thepage}
\renewcommand{\headrulewidth}{0.4pt}

\newcolumntype{Y}{>{\raggedright\arraybackslash}X}
\newcommand{\lp}[1]{\path{#1}}
\newcommand{\code}[1]{\texttt{\small #1}}

\newtcolorbox{keybox}[1][]{
  colback=blue!4, colframe=sgblue!70, boxrule=0.6pt, arc=2pt,
  left=6pt, right=6pt, top=5pt, bottom=5pt, breakable, #1}
\newtcolorbox{warnbox}[1][]{
  colback=red!3, colframe=sgred!60, boxrule=0.6pt, arc=2pt,
  left=6pt, right=6pt, top=5pt, bottom=5pt, breakable, #1}
\newtcolorbox{specbox}[1][]{
  colback=green!3, colframe=sggreen!55, boxrule=0.6pt, arc=2pt,
  left=6pt, right=6pt, top=5pt, bottom=5pt, breakable, #1}

\tikzset{
  onl/.style={draw,rounded corners=2pt,align=center,inner sep=4pt,
              fill=blue!6,font=\footnotesize,minimum height=7mm},
  offl/.style={draw,rounded corners=2pt,align=center,inner sep=4pt,
               fill=green!7,font=\footnotesize,minimum height=7mm},
  store/.style={draw,rounded corners=2pt,align=center,inner sep=5pt,
                fill=orange!12,font=\footnotesize},
  arr/.style={-{Latex[length=1.8mm]},thick,draw=sggray},
}

\title{\bfseries SG-Lean 思路文档（第二版）\\[2mm]
       \large 复用引导的双尺度 Lean 4 证明器\\[1mm]
       \normalsize 外部评审后实现规格}
\author{}
\date{2026 年 9 月 27 日}

\begin{document}
\maketitle
\thispagestyle{fancy}

\begin{keybox}
\textbf{文件地位}：本文件是当前实现规格。历史阶段日志只记录过去发生了什么，
不能覆盖本规格。字段协议见 \lp{sgs-reap/docs/api-contract.md}，数据纪律见
\lp{sgs-reap/docs/data-protocol.md}，工程坑见 \lp{sgs-reap/docs/pitfalls.md}。

\textbf{本次重定稿的核心}：保留“在线单题求解 + 离线建库”两个时间尺度，
但把二者的连接收紧为一个冻结的活动库快照；区分“引理被使用”与“引理真正带来增益”；
删除不成立的子模保证；将现有 35 条库降级为开发期资产。
\end{keybox}

\tableofcontents

%======================================================================
\section{目标、边界与评判指标}

\subsection{产品目标}

输入一条闭式 Lean 4 命题，输出一段通过 Lean 内核终检的 tactic 脚本。

候选必须依次通过 \code{Verify.verify} 和最终 \code{checkProof}；否则不得进入输出字段。

项目不做自然语言自动形式化，不训练或微调模型，不训练检索器，也不重造一个通用的
Mathlib 前提选择系统。研究问题是：\textbf{在零训练条件下，离线构建的可复用引理库，
能否以可控成本提高未参与建库目标的证明成功率？}

\subsection{与 SGS 的关系}

本项目受 SGS（\emph{Scaling Self-Play with Self-Guidance}）的 Solver、Conjecturer、Guide
角色分工启发，但不是 SGS 的训练复现。原 SGS 让合成题通过训练更新 Solver；本项目的
Solver 参数始终冻结，离线积累只通过引理库传递给在线求解器。

因此不能直接继承 SGS 的因果叙事。必须单独验证以下链条：

\[
\text{离线候选}
\rightarrow \text{通过硬门和内核}
\rightarrow \text{跨目标真实使用}
\rightarrow \text{活动库}
\rightarrow \text{在线 pass@}k\text{ 增益}.
\]

\subsection{指标与裁决}

\begin{table}[htbp]
\centering
\small
\begin{tabularx}{\textwidth}{@{}p{2.2cm}p{4.0cm}Y@{}}
\toprule
\textbf{优先级} & \textbf{指标} & \textbf{口径} \\
\midrule
1 & 严格 pass@1 / pass@$k$ & 目标级；首轮模型生成；无 repair、无廉价 tactic \\
2 & 成本 & CostPerSolved、CostPerLemma、CostPerReusable \\
3 & 轻量化 & 0 可训练参数、0 GPU、单实验不超过 12 h、API 不超过 500 元 \\
\bottomrule
\end{tabularx}
\end{table}

廉价 tactic 与 repair 是产品增强，不是库方法本身。它们可以启用，但必须另报
\code{product-pass@k}，不得覆盖严格口径。

\[
\mathrm{CostPerReusable}=
\frac{\text{建库总 token}}{\#\{\ell:\mathrm{reuse}(\ell)\ge 2\}}.
\]

正确率差异显著时取正确率高者；置信区间重叠时取 CostPerSolved 更低者；
成本接近时取实现更简单者。违反轻量化约束的方案直接淘汰。

%======================================================================
\section{固定架构：两个尺度，一个接口}

\subsection{总览}

\begin{center}
\begin{tikzpicture}[node distance=4mm and 5mm]
  \node[onl] (i1) {输入命题};
  \node[onl,right=of i1] (i2) {门检};
  \node[onl,right=of i2] (i3) {相关性召回};
  \node[onl,right=of i3] (i4) {生成 $k$ 篇};
  \node[onl,right=of i4] (i5) {内核终检};
  \node[onl,right=of i5] (i6) {输出通过项};
  \draw[arr] (i1)--(i2); \draw[arr] (i2)--(i3); \draw[arr] (i3)--(i4);
  \draw[arr] (i4)--(i5); \draw[arr] (i5)--(i6);

  \node[store,below=10mm of i3] (snap) {冻结的 active\\LibrarySnapshot};
  \draw[arr] (snap)--(i3);

  \node[offl,below=10mm of snap,xshift=-42mm] (o1) {C-build\\轨迹与需求};
  \node[offl,right=of o1] (o2) {候选与硬门};
  \node[offl,right=of o2] (o3) {probation};
  \node[offl,right=of o3] (o4) {C-measure\\公平曝光};
  \node[offl,right=of o4] (o5) {reuse/cost\\晋升};
  \draw[arr] (o1)--(o2); \draw[arr] (o2)--(o3); \draw[arr] (o3)--(o4);
  \draw[arr] (o4)--(o5); \draw[arr] (o5)--(snap);
\end{tikzpicture}
\end{center}

\begin{keybox}
两个尺度只通过 \textbf{LibrarySnapshot} 相连。在线路径只读，不允许回写全局库；
离线发布新快照后，必须记录快照哈希并重新做环境预检。正式评测期间快照冻结。
\end{keybox}

\subsection{为什么要有试用池}

新引理在产生时必然没有复用证据。如果要求 \code{reuse>0} 才接纳，库无法冷启动；
如果一经证明就进入正式库，活动库又会退化成“所有可证候选的堆积”。因此库必须三态化：

\begin{table}[htbp]
\centering
\small
\begin{tabularx}{\textwidth}{@{}p{2.5cm}Yp{3.2cm}@{}}
\toprule
\textbf{状态} & \textbf{含义} & \textbf{在线可见性} \\
\midrule
probation & 已通过硬门并有内核证明，等待跨目标试用 & 正式在线不可见 \\
active & 在非来源目标上获得足够曝光并达到复用门槛 & 进入冻结快照 \\
cold & 曝光充分仍无复用证据，或从 active 淘汰 & 不可见，保留审计 \\
\bottomrule
\end{tabularx}
\end{table}

%======================================================================
\section{在线单题求解规格}

\subsection{正式方法比较的核心路径}

\begin{enumerate}[leftmargin=1.8em,itemsep=4pt]
  \item \textbf{解析}：把声明或命题字符串转换为闭式命题；失败属于输入错误。
  \item \textbf{门检}：确认能够 elaborate 且类型为 \code{Prop}；协议错误不能伪装成门检拒绝。
  \item \textbf{读取快照}：校验库哈希、物化模块与 import 预检。
  \item \textbf{相关性召回}：先按当前目标召回相关 active 引理。
  \item \textbf{复用排序}：只在相关集合内按 \code{reuse/cost} 排序并做 token 截断。
  \item \textbf{生成}：固定模型、温度、$k$ 与 token 上限，生成 $k$ 篇 tactic 脚本。
  \item \textbf{终检}：批量执行，检查无 sorry、无残余元变量、无未闭合目标，并交内核复核。
  \item \textbf{输出和记账}：只保存通过项；记录快照哈希、候选数、严格 pass@$k$ 与成本。
\end{enumerate}

外部 Mathlib 检索暂不进入核心流程。Mathlib 仍作为 Lean 环境和模型可调用的基础库；
只有在出现稳定端点、明确成本和独立消融之后，才把外部检索作为扩展重新加入。

\subsection{产品增强层}

廉价 tactic 可以在模型前运行，repair 可以在首轮全部失败后运行。二者必须满足：

\begin{itemize}[leftmargin=1.8em,itemsep=3pt]
  \item 最终仍经过同一内核验证器；
  \item 与核心路径共用同一执行引擎，不另写一套验证逻辑；
  \item 报告中分开给出严格口径、cheap 增益、repair 增益和总成本；
  \item 库 A/B 实验默认关闭，以免混淆库效果。
\end{itemize}

\subsection{提示词分离}

\textbf{产品提示词}可以建议优先检查库引理，以提高实际求解率。
\textbf{测量提示词}只能中性列出“这些引理可用”，不得写“优先引用”“更便宜”或
“更不容易错”。否则引用率混入提示服从性，不能解释为复用证据。

%======================================================================
\section{离线建库规格}

\subsection{C 的内部划分}

miniF2F valid 按稳定哈希排序确定性划分为 C-build、C-measure 和 D：

\begin{itemize}[leftmargin=1.8em,itemsep=3pt]
  \item \textbf{C-build}：基线求解、需求挖掘、候选生成和候选证明；
  \item \textbf{C-measure}：给 probation 引理提供跨目标曝光，测量 reuse；
  \item 候选的 \code{source\_target} 无论属于哪一块，都不得为自身贡献 reuse。
\end{itemize}

三个分区精确为 122/61/61，完整覆盖 valid 且两两不相交。分区文件与 manifest
一同入库；manifest 必须记录划分算法、种子、指纹和规模。

\subsection{一轮闭环}

\begin{verbatim}
BUILD_ROUND(C_build, C_measure, snapshot, config):
  1. solve C_build with the current frozen snapshot; collect verified traces
  2. mine demands from normalized subgoals, deduplicated by target
  3. conjecture candidates for unsolved C_build targets
  4. hard-gate candidates: Prop, nontrivial, novel, provable
  5. put verified candidates into probation, not active
  6. schedule neutral, salted exposures on relevant C_measure targets
  7. count constants only from kernel-accepted proofs; exclude source targets
  8. promote by reuse/cost; move fully-tested failures to cold
  9. materialize active entries, compile, preflight, freeze snapshot
 10. persist round report and immutable snapshot manifest
\end{verbatim}

\subsection{硬门}

\begin{table}[htbp]
\centering
\small
\begin{tabularx}{\textwidth}{@{}p{2.0cm}p{3.2cm}Y@{}}
\toprule
\textbf{门} & \textbf{判定} & \textbf{反向对照} \\
\midrule
合式 & elaborate 后是 \code{Prop} & 非 Prop 与缺失响应必须分别失败 \\
非平凡 & 配置内的廉价 tactic 不能证明 & \code{P=P}、\code{x=x}、\code{True} 必须被拒 \\
新颖 & 不与父目标或已有候选/活动库等价 & 改名后的重复命题必须被拒 \\
可证 & 至少一篇证明通过内核终检 & 含 sorry、残余目标或错误常量必须被拒 \\
\bottomrule
\end{tabularx}
\end{table}

证明步数、压缩长度和结构复杂度可以作为诊断字段，但不再成为第二套准入逻辑。

\subsection{复用判据}

对 probation 或 active 引理 $\ell$：

\[
\mathrm{reuse}(\ell)=
\left|\left\{w:\ w\ne\mathrm{source}(\ell),\
\exists y,\ \mathrm{verify}(w,y)=1\land
\ell\in\mathrm{constants}(y)\right\}\right|.
\]

\[
\mathrm{cost}(\ell)=\mathrm{tokens}(\mathrm{render}(\ell)),\qquad
\mathrm{score}(\ell)=\frac{\mathrm{reuse}(\ell)}{\mathrm{cost}(\ell)}.
\]

必须按不同目标去重；失败证明、文本字符串命中和同一目标多次引用都不计数。
每条 probation 引理至少获得规定数量的\textbf{相关目标曝光}后才能判定。曝光调度必须
带轮次盐，防止后端缓存把重复试验变成同一输出的重放。

\begin{warnbox}
\textbf{解释边界}：reuse 是“被使用次数”，不是因果复用价值。它可用于晋升和排序，
但不能单独证明引理提高了解题率。真实价值由冻结的 with/without 配对实验测量。
本项目不再声称随机模型的真实 pass@$k$ 目标具有子模近似保证。
\end{warnbox}

\subsection{晋升、淘汰与快照}

\begin{itemize}[leftmargin=1.8em,itemsep=3pt]
  \item \textbf{晋升}：达到最小曝光数且 \code{reuse >= threshold}，再按 score 和库预算取舍。
  \item \textbf{继续试用}：曝光不足，不能因 reuse 为 0 被淘汰。
  \item \textbf{冷存}：曝光充分仍为 0，或活动库预算下被替换；记录原因与轮次。
  \item \textbf{发布}：只物化 active；编译指定模块；预检成功后生成不可变 snapshot manifest。
\end{itemize}

快照至少记录版本、内容哈希、生成提交、active 条目列表、Lean import、模型配置、
数据指纹和晋升阈值。正式评测报告只引用快照，不引用“当前可变 JSONL”。

%======================================================================
\section{数据协议}

\begin{table}[htbp]
\centering
\small
\begin{tabularx}{\textwidth}{@{}p{1.2cm}p{4.6cm}Y@{}}
\toprule
\textbf{数据} & \textbf{用途} & \textbf{禁止} \\
\midrule
C-build & miniF2F valid 的 122 题；需求和候选 & 晋升测量 \\
C-measure & miniF2F valid 的 61 题；曝光和 reuse & 需求挖掘、候选来源 \\
D & miniF2F valid 的 61 题；调参、debug、难度和预算选择 & 入库、需求挖掘、候选来源 \\
T & miniF2F test；冻结后一次性最终评测 & 任何调参、建库或试跑 \\
\bottomrule
\end{tabularx}
\end{table}

判断泄漏以来源为准，不以文本相似度为准。库条目必须带 \code{source\_target} 与
\code{source\_corpus}。评测时在线产生的临时内容不得回写全局库。

项目的唯一题目来源为 miniF2F。valid 244 题按固定种子
code{sg-lean-minif2f-v1} 和归一化命题哈希切分；test 244 题原样作为 T，
在框架、库快照和统计方案完全冻结前不运行。

%======================================================================
\section{实验设计}

\subsection{三组对照}

\begin{table}[htbp]
\centering
\small
\begin{tabularx}{\textwidth}{@{}p{1.2cm}p{3.4cm}p{3.4cm}Y@{}}
\toprule
\textbf{组} & \textbf{候选生成条件} & \textbf{选择判据} & \textbf{回答的问题} \\
\midrule
A & 真实需求 & SGS Guide & Guide 与形式化复用选择谁更好 \\
B & 真实需求 & reuse/cost & 完整方法 \\
C & 随机等量伪需求 & reuse/cost & 真实需求是否优于随机背景 \\
\bottomrule
\end{tabularx}
\end{table}

A/B 必须使用同一候选池，只改变选择判据；B/C 保持选择、预算和执行配置相同，
只改变需求信号。三组共享模型、温度、$k$、token 上限、硬门、曝光计划、在线检索和验证器。

\subsection{两类测量不能混}

\begin{enumerate}[leftmargin=1.8em,itemsep=4pt]
  \item \textbf{使用测量}：从通过证明的 constants 统计 reuse，用于库生命周期。
  \item \textbf{效果测量}：同一目标、同一配置、冻结随机化方案下做无库/有库配对，
        测严格 pass@$k$ 差值，用于结论。
\end{enumerate}

报告必须给出配对差值、置信区间、按难度分层的结果、后端和协议错误、token 成本，
并抽样做逐引理 with/without 消融。引用代理与真实边际增益的 Spearman 只能作为审计，
不能替代主结果。

\subsection{T 的一次性运行}

只有同时满足以下条件才允许运行 T：代码、超参、C 数据、库快照、评测脚本、随机种子、
报告 schema 全部冻结；D 上装置检查通过；测试台账为空。结果无论好坏都如实报告。

%======================================================================
\section{实现边界与精简原则}

\subsection{一个引擎，三个调用者}

在线 Prover、离线 Runner 和两臂实验必须调用同一个“生成 → Lean 验证 → constants 提取”
执行引擎。它们可以有不同配置，不能复制 HTTP、检索、验证和引用统计实现。

顶层脚本入口仍保持五个：单题、批量评测、离线闭环、配对测量、唯一测试入口。
“一个概念一条实现路径”不等于把所有代码塞进一个大文件；可以拆内部模块，但必须有一个
公共 facade 和一个字段协议。

\subsection{保留、降级与删除}

\begin{table}[htbp]
\centering
\small
\begin{tabularx}{\textwidth}{@{}p{2.0cm}Y@{}}
\toprule
\textbf{类别} & \textbf{内容} \\
\midrule
保留 & Gate、Verify、Trace.constants、Materialize、常驻服务、数据守卫、Guide 对照 \\
核心重构 & 统一执行引擎、库三态、冻结快照、正确两臂统计、相关性优先检索 \\
降级扩展 & cheap、repair、外部 Mathlib 检索；不进入库主实验 \\
退出主线 & Compression/delta\_len、未实测的子模保证、未使用配置与重复路径 \\
\bottomrule
\end{tabularx}
\end{table}

历史日志和 archive 保留事实，但 README、AGENTS、思路文档和接口协议不得继续引用旧结论。

%======================================================================
\section{当前事实与迁移状态}

\subsection{可继续使用的基础}

Gate、内核终检、证明项依赖抽取、物化编译、常驻 Lean 服务、数据角色守卫和逐题落盘
已经具备。纯 Python 自检当前为 95 条断言。

\subsection{旧资产的正确定位}

\begin{warnbox}
现有 \lp{sgs-reap/experiments/library.jsonl} 的 35 条全部 reuse=0，并含 \code{a=a}、
集合基数自反等可由 \code{rfl} 证明的平凡条目。它是开发期装置资产，不是正式活动库。

\code{g3\_c\_device\_n12\_k2.json} 所报“14/24 引用”按生成文本是否包含
\code{sgs\_lem} 统计，包含失败证明，不能作为正式 reuse 证据。

\code{p1\_dev\_k4\_n20.json} 仍是 9 题、\code{partial=true} 的开发报告。
\end{warnbox}

\subsection{代码迁移清单}

\begin{enumerate}[leftmargin=1.8em,itemsep=3pt]
  \item 统一在线、离线和两臂执行路径；
  \item 两臂引用改为通过证明的 Lean constants；
  \item 实现 probation/active/cold 和 snapshot manifest；
  \item 修复非平凡硬门并加入 \code{rfl} 反例；
  \item 改为相关性先召回、reuse/cost 后排序；
  \item 分开产品提示词与中性测量提示词；
  \item 在 miniF2F C-build/C-measure 上重建库，再跑 D 装置检查。
\end{enumerate}

迁移完成前不运行正式 P3，不运行 T。

%======================================================================
\section{阶段计划与验收}

\subsection{P0：连接正确性重构}

交付：统一执行引擎、正确引用统计、库三态、快照协议、硬门反例测试。
闸门：一条失败证明即使写了库名也不计 reuse；一条 \code{rfl} 平凡候选不能进入 probation；
同一来源目标不能为自己的候选贡献 reuse。

\subsection{P1：语料与正式建库}

交付：冻结 miniF2F valid 的 C-build/C-measure/D 分区；在 C 上重建 A/B/C 三个库。
闸门：每个 active 条目均有可追踪的非来源目标引用证据和曝光记录；快照可物化、编译、预检。

\subsection{P2：开发集冻结}

交付：D 上确定模型、$k$、预算、曝光数、阈值、难度层和统计方案。
闸门：严格口径报告可复现；后端错误和协议错误不进入失败分母；配置写入冻结清单。

\subsection{P3：正式对照}

先在 D 上运行 A/B/C；框架完全冻结后，在 T 上一次性运行。任何临时补丁都使 T 结果失效，
必须如实记录而不能“修好再看一次”。

\subsection{P4：成本、忠实度与写作}

产出严格 pass@$k$ 主表、成本 Pareto、reuse 分布、曝光与晋升漏斗、代理忠实度审计、
负结果和限制。论文不得把引用次数写成因果价值，也不得声称不适用的理论保证。

%======================================================================
\appendix
\section{核心字段契约}

\subsection{LibraryEntry}

\begin{verbatim}
{
  "name": "sgs_lem_<content-hash>",
  "stmt": "...",
  "proof": "...",
  "verified": true,
  "source_corpus": "MF_VALID_C",
  "source_target": "aime_...",
  "status": "probation|active|cold",
  "reuse": 3,
  "reuse_targets": ["..."],
  "exposures": 8,
  "exposure_targets": ["..."],
  "cost_tokens": 24,
  "added_round": 2,
  "promoted_round": 4,
  "evicted_round": null
}
\end{verbatim}

\subsection{LibrarySnapshot}

\begin{verbatim}
{
  "version": "snapshot-...",
  "sha256": "...",
  "commit": "...",
  "generated_at": "...",
  "imports": ["Mathlib", "SgsLean.GeneratedLibrary"],
  "active_names": ["sgs_lem_..."],
  "corpus_hashes": {"C_build": "...", "C_measure": "..."},
  "config": {"reuse_threshold": 2, "min_exposures": 4},
  "build_ok": true,
  "preflight_ok": true
}
\end{verbatim}

\subsection{严格评测记录}

\begin{verbatim}
{
  "target": "...",
  "snapshot_hash": "...",
  "k": 4,
  "strict": true,
  "cheap_enabled": false,
  "repair_rounds": 0,
  "proofs_generated": 4,
  "proofs_verified": 1,
  "solved": true,
  "verified_library_constants": ["sgs_lem_..."],
  "usage": {"prompt_tokens": 0, "completion_tokens": 0},
  "backend_error": null,
  "protocol_error": null
}
\end{verbatim}

\end{document}
